\documentclass[aspectratio=169,11pt]{beamer}
\usepackage[T1]{fontenc}
\usepackage[utf8]{inputenc}
\usepackage[french]{babel}
\usepackage{lmodern}
\usepackage{amsmath,amssymb}
\usepackage{graphicx}
\usepackage{booktabs}
\usepackage{tabularx}
\usepackage{xcolor}
\usepackage{tikz}

\usetheme{default}
\definecolor{academicblue}{RGB}{28,61,88}
\setbeamercolor{structure}{fg=academicblue}
\setbeamercolor{normal text}{fg=black,bg=white}
\setbeamercolor{frametitle}{fg=academicblue}
\setbeamerfont{frametitle}{size=\large,series=\bfseries}
\setbeamerfont{title}{size=\Large,series=\bfseries}
\setbeamerfont{subtitle}{size=\normalsize}
\setbeamertemplate{navigation symbols}{}
\setbeamertemplate{footline}{%
  \hspace{0.6cm}\scriptsize Mhamed Halhoul\hfill
  \insertframenumber\hspace{0.6cm}\vspace{0.2cm}}
\setbeamersize{text margin left=0.7cm,text margin right=0.7cm}
\setlength{\emergencystretch}{1em}
\hypersetup{colorlinks=true,urlcolor=academicblue,citecolor=academicblue}
\newcommand{\utrue}{u_{\mathrm{true}}}
\newcommand{\figpath}{../results/extended/figures/}
\newcommand{\figcaption}[1]{\par\vspace{0.1cm}{\footnotesize #1\par}}
\newcommand{\smallgap}{\vspace{0.18cm}}

% Presentation-only legends: same series, colours and line/marker encodings
% as the archived figures. The scientific PDF figures are never rewritten.
\definecolor{legendExpA}{rgb}{0.579607843,0.770196078,0.873725490}
\definecolor{legendExpB}{rgb}{0.331334102,0.622068435,0.804767397}
\definecolor{legendExpC}{rgb}{0.146666667,0.460392157,0.718692810}
\definecolor{legendExpD}{rgb}{0.031372549,0.289734717,0.570319108}
\definecolor{legendRatA}{rgb}{0.992156863,0.656470588,0.382745098}
\definecolor{legendRatB}{rgb}{0.966320646,0.475432526,0.147020377}
\definecolor{legendRatC}{rgb}{0.863529412,0.299607843,0.013333333}
\definecolor{legendRatD}{rgb}{0.617993080,0.199077278,0.012610534}
\definecolor{legendStepA}{rgb}{0.273006000,0.204520000,0.501721000}
\definecolor{legendStepB}{rgb}{0.127568000,0.566949000,0.550556000}
\definecolor{legendStepC}{rgb}{0.606045000,0.850733000,0.236712000}
\definecolor{legendNoisy}{gray}{0.45}
\newcommand{\legendLine}[3]{%
  \begin{tikzpicture}[baseline=-0.6ex]
    \draw[color=#1,#2,line width=0.8pt] (0,0)--(0.44,0);
    \ifcase#3\relax
    \or\fill[color=#1] (0.22,0) circle (1.4pt);
    \or\fill[color=#1] (0.18,-0.04) rectangle (0.26,0.04);
    \fi
  \end{tikzpicture}}
\newcommand{\metricLegend}{%
  {\footnotesize\renewcommand{\arraystretch}{1.08}%
  \begin{tabular}{@{}l@{\hspace{0.35em}}l@{}}
    \legendLine{black}{solid}{1} & Chaleur \\
    \legendLine{legendExpA}{solid}{1} & PM exp., $K=0.025$ \\
    \legendLine{legendExpB}{solid}{1} & PM exp., $K=0.05$ \\
    \legendLine{legendExpC}{solid}{1} & PM exp., $K=0.1$ \\
    \legendLine{legendExpD}{solid}{1} & PM exp., $K=0.2$ \\
    \legendLine{legendRatA}{dashed}{1} & PM rat., $K=0.025$ \\
    \legendLine{legendRatB}{dashed}{1} & PM rat., $K=0.05$ \\
    \legendLine{legendRatC}{dashed}{1} & PM rat., $K=0.1$ \\
    \legendLine{legendRatD}{dashed}{1} & PM rat., $K=0.2$ \\
    \legendLine{legendNoisy}{densely dotted}{0} & Bruit\'ee (r\'ef.) \\
  \end{tabular}}}
\newcommand{\sensitivityLegend}{%
  {\footnotesize\renewcommand{\arraystretch}{1.08}%
  \begin{tabular}{@{}l@{\hspace{0.35em}}l@{}}
    \legendLine{legendStepA}{densely dotted}{0} & Chaleur, $n=5$ \\
    \legendLine{legendStepB}{densely dotted}{0} & Chaleur, $n=20$ \\
    \legendLine{legendStepC}{densely dotted}{0} & Chaleur, $n=80$ \\
    \legendLine{legendStepA}{solid}{1} & PM exp., $n=5$ \\
    \legendLine{legendStepA}{dashed}{2} & PM rat., $n=5$ \\
    \legendLine{legendStepB}{solid}{1} & PM exp., $n=20$ \\
    \legendLine{legendStepB}{dashed}{2} & PM rat., $n=20$ \\
    \legendLine{legendStepC}{solid}{1} & PM exp., $n=80$ \\
    \legendLine{legendStepC}{dashed}{2} & PM rat., $n=80$ \\
    \legendLine{legendNoisy}{densely dotted}{0} & Bruit\'ee (r\'ef.) \\
  \end{tabular}}}

\title{Traitement d'images par des équations aux dérivées partielles}
\subtitle{Débruitage par diffusion linéaire\\et diffusion de Perona--Malik}
\author{Mhamed Halhoul}
\institute{M2 -- Modélisation, Mathématiques Appliquées et Apprentissage\\
Abdelmalek Essaâdi University}
\date{}

\begin{document}

% 25 diapositives principales. Les labels correspondent aux notes orales.
\begin{frame}[label=titre,plain]
  \titlepage
\end{frame}

\begin{frame}[label=question]{Réduction du bruit et conservation des structures}
  \large
  Dans quelles conditions la diffusion dépendant des variations locales
  modifie-t-elle le compromis entre réduction du bruit et préservation
  des contours, par rapport à la chaleur ?
  \smallgap
  \normalsize
  \begin{itemize}
    \item Une comparaison de modèles et de schémas numériques.
    \item Une étude reproductible sur deux images synthétiques.
    \item Aucune supériorité de Perona--Malik posée à l'avance.
  \end{itemize}
\end{frame}

\begin{frame}[label=observation]{Image propre et observation bruitée}
  \[
    \utrue:\Omega\longrightarrow[0,1],\qquad
    f_{ij}=(\utrue)_{ij}+\sigma Z_{ij},\quad
    Z_{ij}\overset{\mathrm{iid}}{\sim}\mathcal N(0,1).
  \]
  \begin{itemize}
    \item Rectangle spatial, puis matrice de $128\times128$ pixels.
    \item Une observation commune à toutes les méthodes pour chaque cas.
    \item Calculs non tronqués : $f$ et les sorties peuvent dépasser $[0,1]$.
  \end{itemize}
  \smallgap
  L'intervalle $[0,1]$ fixe l'échelle d'affichage des images.
  Les profils et les métriques conservent les valeurs numériques.
\end{frame}

\begin{frame}[label=chaleur]{La chaleur diffuse toutes les variations}
  \[
    \partial_tu=\Delta u,\qquad u(\cdot,0)=f,\qquad
    \partial_nu=0\ \text{sur }\partial\Omega.
  \]
  \begin{itemize}
    \item La diffusion de chaleur atténue surtout les hautes fréquences.
    \item Le flux nul conserve la moyenne.
    \item Une diffusion prolongée élargit aussi les transitions utiles.
  \end{itemize}
  \smallgap
  Sur $\mathbb R^2$, convolution par une gaussienne de variance $2t$.
  Sur le rectangle, le bord de Neumann change le noyau de solution.
\end{frame}

\begin{frame}[label=pm]{Une conductance dépend de la variation locale}
  \[
    \partial_tu=\operatorname{div}\bigl(c(|\nabla u|)\nabla u\bigr),
    \qquad u(\cdot,0)=f.
  \]
  \[
    c_{\exp}(s)=e^{-(s/K)^2},\qquad
    c_{\mathrm{rat}}(s)=\frac{1}{1+(s/K)^2},\qquad K>0.
  \]
  \begin{itemize}
    \item Une forte variation réduit la conductance.
    \item Un faible $K$ peut préserver les contours \emph{et} le bruit.
    \item Le résultat dépend ensemble de $K$, de la conduction et du temps.
  \end{itemize}
  {\footnotesize Modèle de Perona et Malik \cite{perona1990}.
  Le flux $s\,c(s)$ peut décroître : difficulté du modèle continu.}
\end{frame}

\begin{frame}[label=directionnel]{Le solveur emploie une variante directionnelle}
  Entre deux pixels horizontaux :
  \[
    \delta_x=u^n_{i,j+1}-u^n_{ij},\qquad
    a^n_{i,j+1/2}=c\!\left(\frac{|\delta_x|}{dx}\right).
  \]
  \begin{itemize}
    \item Quatre voisins : deux différences horizontales et deux verticales.
    \item La conductance utilise une différence directionnelle par arête.
    \item Aucune reconstruction de la norme complète $|\nabla u|$ à l'interface.
  \end{itemize}
  \smallgap
  Variante discrète directionnelle à quatre voisins, sensible à
  l'orientation de la grille.
\end{frame}

\begin{frame}[label=flux]{Une arête intérieure partage un flux opposé}
  Contribution à la divergence de l'arête $(i,j)$--$(i,j+1)$ :
  \[
    F_x= a^n_{i,j+1/2}\,\frac{u^n_{i,j+1}-u^n_{ij}}{dx^2}.
  \]
  \[
    (L_hu^n)_{ij}\mathrel{+}=F_x,
    \qquad (L_hu^n)_{i,j+1}\mathrel{-}=F_x.
  \]
  \begin{itemize}
    \item Même conductance pour les deux pixels de l'arête.
    \item Aucune arête extérieure : flux nul, sans connexion périodique.
    \item Euler explicite : $u^{n+1}=u^n+dt\,L_h(u^n)$.
  \end{itemize}
  \smallgap
  Tous les flux partent de $u^n$. Les checkpoints conservent des copies
  d'états, sans mélange de deux niveaux temporels.
\end{frame}

\begin{frame}[label=stabilite]{Une condition suffisante donne le maximum discret}
  \[
    dt\left(\frac{1}{dx^2}+\frac{1}{dy^2}\right)\leq\frac12,
    \qquad dx=dy=1,\quad dt=0{,}20.
  \]
  \[
    u^{n+1}_p=
    \left(1-\sum_{q\sim p}\alpha_{pq}^n\right)u^n_p
      +\sum_{q\sim p}\alpha_{pq}^n u^n_q,
    \qquad \alpha_{pq}^n=\frac{dt\,a_{pq}^n}{h_{pq}^2}.
  \]
  \begin{itemize}
    \item $0\leq a_{pq}^n\leq1$ : coefficients non négatifs, somme égale à un.
    \item Constantes conservées et valeurs entre les extrema initiaux.
    \item Flux opposés : conservation de la somme et de la moyenne.
  \end{itemize}
  {\footnotesize Ces propriétés du schéma ne prouvent pas la bonne position
  du modèle continu de Perona--Malik.}
\end{frame}

\begin{frame}[label=protocole]{Une grille fixée après l'exploration rapide}
  \begin{columns}[T,onlytextwidth]
    \column{0.53\textwidth}
    \begin{itemize}
      \item Deux images synthétiques de $128\times128$.
      \item $\sigma\in\{0{,}025;0{,}05;0{,}10\}$.
      \item Dix graines : $0,\ldots,9$.
      \item Chaleur et huit variantes PM.
    \end{itemize}
    \column{0.44\textwidth}
    \begin{itemize}
      \item $K\in\{0{,}025;0{,}05;0{,}10;0{,}20\}$.
      \item $n\in\{0,1,2,5,10,20,40,80\}$.
      \item Temps de diffusion $t=n\,dt$.
      \item $60$ observations, $540$ trajectoires.
    \end{itemize}
  \end{columns}
  \smallgap
  $4\,380$ lignes de mesures et $438$ groupes.
  Dix réalisations de bruit par image, et non dix images.
  Une graine recycle le même champ standardisé entre images et $\sigma$.
\end{frame}

\begin{frame}[label=metriques]{PSNR et SSIM décrivent des aspects différents}
  \[
    \mathrm{MSE}=\frac{\|u-\utrue\|_2^2}{N},\quad N=N_xN_y,\qquad
    \mathrm{PSNR}=10\log_{10}\!\left(\frac{1}{\mathrm{MSE}}\right).
  \]
  \begin{itemize}
    \item Erreur relative : $\|u-\utrue\|_2/\|\utrue\|_2$.
    \item SSIM : luminance, contraste et structure dans des fenêtres $7\times7$.
    \item Fenêtres uniformes valides, covariances échantillonnales.
    \item $K_1=0{,}01$, $K_2=0{,}03$, \texttt{data\_range=1}.
  \end{itemize}
  \smallgap
  SSIM peut être négatif. Moyennes des scores individuels.
  Bandes et barres : un écart-type échantillonnal, sans intervalle de confiance.
  {\footnotesize SSIM original : Wang et al. \cite{wang2004}.
  Variante documentée : \cite{skimage_metrics}.}
\end{frame}

\begin{frame}[label=verification]{Vérification et reproductibilité ont deux niveaux}
  \begin{itemize}
    \item $165$ tests réussis lors des validations et de l'audit antérieurs.
    \item Chaleur : mode propre discret de Neumann et contrôles de raffinement.
    \item Mode strict : égalité exacte obtenue dans le Windows enregistré.
    \item Audit communiqué : mode numerical inter-environnements,
      $\mathrm{rtol}=10^{-10}$, $\mathrm{atol}=10^{-12}$.
  \end{itemize}
  \smallgap
  Les manifestes contrôlent l'intégrité exacte des fichiers archivés.
  Le mode numerical ne revendique pas une reproduction bit à bit.
  \smallgap
  {\footnotesize Pipeline enregistré : $44{,}407302$ s dans l'environnement
  d'origine. Durées solveurs cumulatives aux checkpoints, distinctes du
  temps de diffusion. Aucun nouveau calcul scientifique pour cet exposé.}
\end{frame}

\begin{frame}[label=comparaison]{Le cas visuel fixé avant calcul}
  \centering
  \includegraphics[width=\textwidth,height=0.50\textheight,keepaspectratio]
    {\figpath comparison.pdf}
  \figcaption{\texttt{geometric\_shapes}, $\sigma=0{,}10$, graine $0$.
  Chaleur et PM : $n=20$, $t=4$, PM : $K=0{,}10$.
  Même observation, affichage $[0,1]$, calculs non tronqués.}
  \smallgap
  \raggedright
  La chaleur élargit les transitions. L'exponentielle laisse des pixels
  isolés. La rationnelle atténue davantage ces fluctuations dans ce cas.
\end{frame}

\begin{frame}[label=profil]{Un profil révèle les pics et la largeur des transitions}
  \centering
  \includegraphics[width=\textwidth,height=0.64\textheight,keepaspectratio]
    {\figpath intensity_profile.pdf}
  \figcaption{Profil horizontal ligne $42$ : \texttt{geometric\_shapes},
  $\sigma=0{,}10$, graine $0$, $n=20$, $K=0{,}10$.
  Intensités non tronquées, sans agrégation.}
\end{frame}

\begin{frame}[label=gradients]{Les cartes de gradient visualisent les structures}
  \centering
  \includegraphics[width=\textwidth,height=0.48\textheight,keepaspectratio]
    {\figpath gradient_maps.pdf}
  \figcaption{\texttt{geometric\_shapes}, $\sigma=0{,}10$, graine $0$,
  $n=20$, $K=0{,}10$. Norme calculée par \texttt{np.gradient},
  même échelle de couleurs, plafond au maximum des 99es percentiles.}
  \smallgap
  \raggedright
  Le bruit résiduel apparaît aussi dans le gradient.
  Cette visualisation n'ajoute aucune métrique quantitative de contours.
\end{frame}

\begin{frame}[label=courbes-formes]{Formes géométriques : PSNR et temps de diffusion}
  % Recadrages du PDF archivé : panneau gauche ; légende recomposée à partir des mêmes séries.
  % trim = gauche, bas, droite, haut, en points PDF (bp).
  \begin{columns}[T,onlytextwidth]
    \column{0.71\textwidth}
    \includegraphics[width=\linewidth,trim=0bp 118bp 462bp 90bp,clip]
      {\figpath metric_curves_geometric_shapes_sigma0p1.pdf}
    \column{0.27\textwidth}
    \centering
    {\small Méthodes et seuils $K$\par}
    \smallgap
    \metricLegend
  \end{columns}
  \figcaption{\texttt{geometric\_shapes}, $\sigma=0{,}10$, dix graines.
  Bandes : $\pm$ un écart-type échantillonnal. $dt=0{,}20$,
  $n\in\{0,1,2,5,10,20,40,80\}$. Toutes les configurations sont visibles.}
\end{frame}

\begin{frame}[label=courbes-formes-ssim]{Formes géométriques : SSIM et temps de diffusion}
  \begin{columns}[T,onlytextwidth]
    \column{0.71\textwidth}
    \includegraphics[width=\linewidth,trim=463bp 118bp 0bp 90bp,clip]
      {\figpath metric_curves_geometric_shapes_sigma0p1.pdf}
    \column{0.27\textwidth}
    \centering
    {\small Méthodes et seuils $K$\par}
    \smallgap
    \metricLegend
  \end{columns}
  \figcaption{\texttt{geometric\_shapes}, $\sigma=0{,}10$, dix graines.
  Bandes : $\pm$ un écart-type échantillonnal. $dt=0{,}20$,
  $n\in\{0,1,2,5,10,20,40,80\}$. Mêmes trajectoires que le PSNR.}
\end{frame}

\begin{frame}[label=courbes-rampes]{Rampes et contours : PSNR et temps de diffusion}
  \begin{columns}[T,onlytextwidth]
    \column{0.71\textwidth}
    \includegraphics[width=\linewidth,trim=0bp 118bp 462bp 90bp,clip]
      {\figpath metric_curves_ramps_and_edges_sigma0p1.pdf}
    \column{0.27\textwidth}
    \centering
    {\small Méthodes et seuils $K$\par}
    \smallgap
    \metricLegend
  \end{columns}
  \figcaption{\texttt{ramps\_and\_edges}, $\sigma=0{,}10$, dix graines.
  Bandes : $\pm$ un écart-type échantillonnal. $dt=0{,}20$,
  $n\in\{0,1,2,5,10,20,40,80\}$. Toutes les configurations sont visibles.}
\end{frame}

\begin{frame}[label=courbes-rampes-ssim]{Rampes et contours : SSIM et temps de diffusion}
  \begin{columns}[T,onlytextwidth]
    \column{0.71\textwidth}
    \includegraphics[width=\linewidth,trim=463bp 118bp 0bp 90bp,clip]
      {\figpath metric_curves_ramps_and_edges_sigma0p1.pdf}
    \column{0.27\textwidth}
    \centering
    {\small Méthodes et seuils $K$\par}
    \smallgap
    \metricLegend
  \end{columns}
  \figcaption{\texttt{ramps\_and\_edges}, $\sigma=0{,}10$, dix graines.
  Bandes : $\pm$ un écart-type échantillonnal. $dt=0{,}20$,
  $n\in\{0,1,2,5,10,20,40,80\}$. Mêmes trajectoires que le PSNR.}
\end{frame}

\begin{frame}[label=sensibilite]{Formes géométriques : sensibilité du PSNR à K}
  \begin{columns}[T,onlytextwidth]
    \column{0.71\textwidth}
    \includegraphics[width=\linewidth,trim=0bp 118bp 462bp 90bp,clip]
      {\figpath k_sensitivity_geometric_shapes_sigma0p1.pdf}
    \column{0.27\textwidth}
    \centering
    {\small Méthodes et checkpoints\par}
    \smallgap
    \sensitivityLegend
  \end{columns}
  \figcaption{\texttt{geometric\_shapes}, $\sigma=0{,}10$, dix graines.
  $n=5,20,80$ ($t=1,4,16$). Barres : $\pm$ un écart-type échantillonnal.
  Tous les $K\in\{0{,}025;0{,}05;0{,}10;0{,}20\}$, axe logarithmique.}
\end{frame}

\begin{frame}[label=sensibilite-ssim]{Formes géométriques : sensibilité du SSIM à K}
  \begin{columns}[T,onlytextwidth]
    \column{0.71\textwidth}
    \includegraphics[width=\linewidth,trim=463bp 118bp 0bp 90bp,clip]
      {\figpath k_sensitivity_geometric_shapes_sigma0p1.pdf}
    \column{0.27\textwidth}
    \centering
    {\small Méthodes et checkpoints\par}
    \smallgap
    \sensitivityLegend
  \end{columns}
  \figcaption{\texttt{geometric\_shapes}, $\sigma=0{,}10$, dix graines.
  $n=5,20,80$ ($t=1,4,16$). Barres : $\pm$ un écart-type échantillonnal.
  Tous les $K\in\{0{,}025;0{,}05;0{,}10;0{,}20\}$, axe logarithmique.}
\end{frame}

\begin{frame}[label=desaccord]{PSNR et SSIM peuvent classer différemment}
  \centering
  {\small\renewcommand{\arraystretch}{1.35}
  % Generated from archived summary.csv; rounding applies only to display.
\begin{tabular}{lll}
\toprule
Méthode & PSNR (dB) & SSIM \\
\midrule
Bruitée & $19.9989 \pm 0.0458$ & $0.2168 \pm 0.0013$ \\
Chaleur & $23.6482 \pm 0.0432$ & $0.8017 \pm 0.0023$ \\
PM exp. & $28.3449 \pm 0.2351$ & $0.6254 \pm 0.0137$ \\
PM rat. & $34.6530 \pm 0.2974$ & $0.9292 \pm 0.0027$ \\
\bottomrule
\end{tabular}
}
  \figcaption{\texttt{geometric\_shapes}, $\sigma=0{,}10$.
  Chaleur et PM : $n=20$ ($t=4$), PM : $K=0{,}10$.
  Bruitée : $n=0$. Moyennes $\pm$ écarts-types sur dix graines.}
  \smallgap
  \raggedright
  L'exponentielle dépasse la chaleur en PSNR, mais son SSIM est inférieur.
  La rationnelle est meilleure selon les deux scores dans cette comparaison
  à paramètres fixés.
\end{frame}

\begin{frame}[label=oracle]{Oracles exploratoires sur la grille étudiée}
  \[
    \theta_{\mathrm{oracle}}
      \in\arg\max_{\theta\in\mathcal G,\ n(\theta)>0}
        \frac1{10}\sum_{r=0}^{9}
          \mathrm{PSNR}(u^{(r)}_{\theta},\utrue).
  \]
  \begin{itemize}
    \item $\mathcal G$ : chaleur et PM, checkpoints et paramètres prédéfinis.
    \item Sélection avec accès à l'image propre, après mesure.
    \item Un oracle SSIM peut choisir un autre réglage.
    \item Un maximum à $n=80$ reste à la frontière de la grille.
  \end{itemize}
  \smallgap
  Cet oracle exploratoire décrit les configurations étudiées.
  Il ne fournit pas une règle de choix de $K$ ou d'arrêt en pratique.
\end{frame}

\begin{frame}[label=limites]{Portée de l'étude et limites}
  \begin{itemize}
    \item Deux images synthétiques et bruit gaussien seulement.
    \item Dix bruits par image, cas liés par la réutilisation des graines.
    \item Grille limitée de paramètres, construite après exploration.
    \item Variante directionnelle, métriques avec accès à la référence.
    \item Aucune validation sur images naturelles, aucun test de supériorité.
  \end{itemize}
  \smallgap
  Perspectives : provenance vérifiée d'images naturelles, sélection sans
  référence, autres schémas et régularisation de Catté et al.
  \cite{catte1992,weickert1998}.
\end{frame}

\begin{frame}[label=conclusion]{Une réponse conditionnelle à la question scientifique}
  \large
  La diffusion dépendant des variations locales change le compromis,
  mais son intérêt dépend du seuil, du temps, de l'image et du critère.
  \smallgap
  \normalsize
  \begin{itemize}
    \item Dans le cas fixé, la rationnelle conserve mieux les structures
      visibles tout en réduisant le bruit.
    \item Un faible $K$ peut préserver du bruit, un temps long peut sur-lisser.
    \item Un gain en PSNR n'impose pas un gain en SSIM.
  \end{itemize}
  \smallgap
  Ces constats portent sur la grille mesurée et les deux images synthétiques.
\end{frame}

\begin{frame}[label=references]{Références mobilisées}
  \footnotesize
  \setbeamertemplate{bibliography item}[text]
  \begin{thebibliography}{5}
    \bibitem{perona1990} P. Perona et J. Malik.
      \emph{Scale-Space and Edge Detection Using Anisotropic Diffusion}.
      IEEE TPAMI, 1990.
      \href{https://doi.org/10.1109/34.56205}{doi:10.1109/34.56205}.
    \bibitem{weickert1998} J. Weickert.
      \emph{Anisotropic Diffusion in Image Processing}.
      Teubner, 1998.
      \href{https://www.mia.uni-saarland.de/weickert/book.html}{Texte sur le site de l'auteur}.
    \bibitem{catte1992} F. Catté, P.-L. Lions, J.-M. Morel et T. Coll.
      \emph{Image Selective Smoothing and Edge Detection by Nonlinear Diffusion}.
      SIAM J. Numer. Anal., 1992.
      \href{https://doi.org/10.1137/0729012}{doi:10.1137/0729012}.
    \bibitem{wang2004} Z. Wang, A. C. Bovik, H. R. Sheikh et E. P. Simoncelli.
      \emph{Image Quality Assessment: From Error Visibility to Structural Similarity}.
      IEEE Trans. Image Process., 2004.
      \href{https://doi.org/10.1109/TIP.2003.819861}{doi:10.1109/TIP.2003.819861}.
    \bibitem{skimage_metrics} Documentation officielle scikit-image.
      \href{https://scikit-image.org/docs/stable/api/skimage.metrics.html}{Métriques PSNR et SSIM}.
  \end{thebibliography}
\end{frame}

\appendix
\section{Réserves}

\begin{frame}[label=reserve-conservation]{Réserve : conservation par annulation des flux}
  \[
    \begin{aligned}
      \sum_p u^{n+1}_p-\sum_p u^n_p
      &=dt\sum_{\{p,q\}\in\mathcal E}
       \left(a_{pq}^n\frac{u_q^n-u_p^n}{h_{pq}^2}
            +a_{pq}^n\frac{u_p^n-u_q^n}{h_{pq}^2}\right)\\
      &=0.
    \end{aligned}
  \]
  \begin{itemize}
    \item Une arête non orientée apparaît une fois, avec deux contributions.
    \item La grille est uniforme et les volumes des cellules sont identiques.
    \item Pas de contribution extérieure au bord.
  \end{itemize}
  \smallgap
  Pour $u^n$ constant, chaque différence est nulle : $u^{n+1}=u^n$.
  Les égalités sont algébriques, à l'arrondi près en virgule flottante.
\end{frame}

\begin{frame}[label=reserve-maximum]{Réserve : coefficients de la combinaison convexe}
  \[
    \sum_{q\sim p}\alpha_{pq}^n
     \leq 2dt\left(\frac1{dx^2}+\frac1{dy^2}\right)\leq1.
  \]
  \[
    m_n\leq u_p^{n+1}\leq M_n,
    \qquad m_n=\min_pu_p^n,\quad M_n=\max_pu_p^n.
  \]
  \begin{itemize}
    \item Au bord, il y a moins de voisins : la même borne suffit.
    \item Pour PM, les poids dépendent de $u^n$, mais restent dans $[0,1]$.
    \item Les extrema de l'observation peuvent être hors de $[0,1]$.
  \end{itemize}
  \smallgap
  Aucune contraction entre deux solutions ni décroissance d'une énergie
  non linéaire n'est déduite de cette seule démonstration.
\end{frame}

\begin{frame}[label=reserve-ssim]{Réserve : définition exacte du SSIM calculé}
  \[
    S_W(x,y)=
    \frac{(2\mu_x\mu_y+C_1)(2s_{xy}+C_2)}
    {(\mu_x^2+\mu_y^2+C_1)(s_x^2+s_y^2+C_2)},
    \qquad C_1=0{,}01^2,\quad C_2=0{,}03^2.
  \]
  \begin{itemize}
    \item Moyennes uniformes sur $49$ pixels, variances avec dénominateur $48$.
    \item Moyenne de $S_W$ sur les seules fenêtres complètes $7\times7$.
    \item Une covariance négative peut produire un SSIM négatif.
    \item Différence avec la fenêtre gaussienne de la publication d'origine.
  \end{itemize}
  \smallgap
  Le PSNR moyen est la moyenne des logarithmes individuels.
  Il diffère en général du PSNR de la MSE moyenne.
\end{frame}

\begin{frame}[label=reserve-neumann]{Réserve : un mode propre discret de Neumann}
  \[
    \phi_{pq}[i,j]=
      \cos\!\left(\frac{\pi p(i+1/2)}{N_y}\right)
      \cos\!\left(\frac{\pi q(j+1/2)}{N_x}\right),
  \]
  \[
    \lambda_{pq}=-\frac4{dy^2}\sin^2\!\left(\frac{\pi p}{2N_y}\right)
                 -\frac4{dx^2}\sin^2\!\left(\frac{\pi q}{2N_x}\right),
    \qquad u^n=(1+dt\,\lambda_{pq})^n\phi_{pq}.
  \]
  Cette amplitude fournit une référence indépendante du solveur.
  \smallgap
  Deux niveaux de raffinement contrôlent séparément une diminution d'erreur
  spatiale et temporelle. Ils ne prouvent pas un ordre général de convergence.
\end{frame}

\begin{frame}[label=reserve-sensibilite-rampes]{Réserve : sensibilité du PSNR sur les rampes}
  \begin{columns}[T,onlytextwidth]
    \column{0.71\textwidth}
    \includegraphics[width=\linewidth,trim=0bp 118bp 462bp 90bp,clip]
      {\figpath k_sensitivity_ramps_and_edges_sigma0p1.pdf}
    \column{0.27\textwidth}
    \centering
    {\small Méthodes et checkpoints\par}
    \smallgap
    \sensitivityLegend
  \end{columns}
  \figcaption{\texttt{ramps\_and\_edges}, $\sigma=0{,}10$, dix graines.
  $n=5,20,80$ ($t=1,4,16$). Barres : $\pm$ un écart-type échantillonnal.
  Tous les $K\in\{0{,}025;0{,}05;0{,}10;0{,}20\}$, axe logarithmique.}
\end{frame}

\begin{frame}[label=reserve-sensibilite-rampes-ssim]{Réserve : sensibilité du SSIM sur les rampes}
  \begin{columns}[T,onlytextwidth]
    \column{0.71\textwidth}
    \includegraphics[width=\linewidth,trim=463bp 118bp 0bp 90bp,clip]
      {\figpath k_sensitivity_ramps_and_edges_sigma0p1.pdf}
    \column{0.27\textwidth}
    \centering
    {\small Méthodes et checkpoints\par}
    \smallgap
    \sensitivityLegend
  \end{columns}
  \figcaption{\texttt{ramps\_and\_edges}, $\sigma=0{,}10$, dix graines.
  $n=5,20,80$ ($t=1,4,16$). Barres : $\pm$ un écart-type échantillonnal.
  Tous les $K\in\{0{,}025;0{,}05;0{,}10;0{,}20\}$, axe logarithmique.}
\end{frame}

\begin{frame}[label=reserve-continu]{Réserve : dérivée du flux et modèle continu}
  Poser $g(s)=s\,c(s)$, $s\geq0$ :
  \[
    g'_{\exp}(s)=e^{-(s/K)^2}\bigl(1-2(s/K)^2\bigr),
    \qquad
    g'_{\mathrm{rat}}(s)=\frac{1-(s/K)^2}{(1+(s/K)^2)^2}.
  \]
  \begin{itemize}
    \item Dérivées négatives pour $s>K/\sqrt2$ et $s>K$, respectivement.
    \item Le modèle continu peut perdre la parabolicité dans la direction
      du gradient.
    \item $0\leq c\leq1$ et la condition de pas donnent le maximum discret,
      sans résoudre cette difficulté continue.
  \end{itemize}
  \smallgap
  Catté et al. évaluent la conductance sur un gradient spatialement lissé.
  Cette régularisation bibliographique n'a pas été implémentée.
\end{frame}

\end{document}
